\chapter{Wymagania systemowe}
\label{cha:wymagania}

\section{Wymagania funkcjonalne}
\label{sec:wymaganiaFunkcjonalne}

\begin{itemize}
  \item \textbf{Zarządzanie kontami i rolami użytkowników}
  \begin{itemize}
    \item Użytkownik może samodzielnie zarejestrować nowe konto (imię, nazwisko, login, hasło),
          wybierając przy rejestracji rolę - ucznia lub nauczyciela. Login jest unikalny w skali
          całego systemu.
    \item System rozróżnia dwie role - nauczyciela i ucznia - i po zalogowaniu udostępnia każdej
          z nich odrębny zestaw funkcji oraz inny układ nawigacji. Każdy nauczyciel zarządza
          wyłącznie własnymi testami.
  \end{itemize}

  \item \textbf{Zarządzanie grupami uczniów}
  \begin{itemize}
    \item Nauczyciel może tworzyć nowe grupy opatrzone nazwą i opisem oraz zarządzać ich składem.
    \item Jeden uczeń może należeć do wielu grup jednocześnie.
  \end{itemize}

  \item \textbf{Tworzenie i edycja testów}
  \begin{itemize}
    \item Nauczyciel tworzy testy złożone z pytań trzech typów: jednokrotnego wyboru, wielokrotnego
          wyboru oraz pytań otwartych.
    \item Dla każdego testu określa nazwę, przedmiot, limit czasu, próg zaliczenia, dozwoloną liczbę
          podejść, widoczność oraz grupy, którym test jest udostępniany.
    \item Test można edytować, duplikować i usuwać. Po pojawieniu się rozwiązań edycja pytań jest
          blokowana, aby nie naruszyć powiązanych z nimi odpowiedzi.
    \item Tylko aktywny test jest widoczny dla uczniów; nowy test powstaje jako nieaktywny (szkic).
  \end{itemize}

  \item \textbf{Rozwiązywanie testów}
  \begin{itemize}
    \item Uczeń widzi listę testów przypisanych do jego grup, z podziałem na rozwiązane
          i nierozwiązane.
    \item Podczas rozwiązywania działa licznik czasu; po jego upływie test jest automatycznie
          oddawany. Liczba podejść jest ograniczona wartością ustawioną przez nauczyciela.
  \end{itemize}

  \item \textbf{Ocenianie odpowiedzi}
  \begin{itemize}
    \item Pytania zamknięte są oceniane automatycznie, w wielokrotnym wyborze z punktami częściowymi.
    \item Pytania otwarte ocenia ręcznie nauczyciel; może też skorygować punkty dowolnego pytania
          każdego podejścia. Przy wielu podejściach liczy się najlepszy wynik.
  \end{itemize}

  \item \textbf{Wykrywanie prób oszustwa}
  \begin{itemize}
    \item Opuszczenie okna testu jest wykrywane: pierwsze skutkuje ostrzeżeniem, kolejne -
          zakończeniem testu z wynikiem 0\% i oznaczeniem próby oszustwa.
  \end{itemize}

  \item \textbf{Statystyki i eksport wyników}
  \begin{itemize}
    \item Nauczyciel ma dostęp do statystyk testu (zdawalność, średni wynik, najtrudniejsze pytanie,
          poprawność pytań) oraz może wyeksportować wyniki do pliku CSV.
  \end{itemize}
\end{itemize}

\section{Wymagania niefunkcjonalne}
\label{sec:wymaganiaNiefunkcjonalne}

\textbf{Bezpieczeństwo}
\begin{itemize}
  \item Dostęp do funkcji systemu jest możliwy wyłącznie po zalogowaniu, a zakres uprawnień zależy
        od roli. Hasła nie są przechowywane jawnie - zapisuje się ich skrót (SHA-256). Unikalność
        loginów jest wymuszana na poziomie bazy danych.
\end{itemize}

\textbf{Niezawodność i obsługa błędów}
\begin{itemize}
  \item Operacje na bazie danych oraz operacje wejścia/wyjścia są zabezpieczone wspólnym mechanizmem
        obsługi wyjątków. W razie błędu użytkownik otrzymuje zrozumiały komunikat, a aplikacja nie
        ulega awarii.
\end{itemize}

\textbf{Spójność danych}
\begin{itemize}
  \item Integralność jest wymuszana kluczami obcymi z odpowiednimi regułami usuwania, unikalnym
        indeksem oraz mechanizmem migracji wersjonującym schemat bazy.
\end{itemize}

\textbf{Użyteczność}
\begin{itemize}
  \item Interfejs jest czytelny i spójny, nawigacja odbywa się przez stały panel boczny, a komunikaty
        prezentowane są w języku polskim. Nowy użytkownik bez trudu odnajduje podstawowe funkcje.
\end{itemize}

\textbf{Wydajność}
\begin{itemize}
  \item Kluczowe operacje (logowanie, wczytywanie testów, zapis wyniku) wykonują się w czasie
        nieprzekraczającym kilku sekund.
\end{itemize}

\section{Założenia dotyczące działania systemu}
\label{sec:zalozeniaSystemu}

Przyjęto następujące założenia:

\begin{itemize}
  \item Aplikacja działa w systemie Windows 10/11 z zainstalowanym środowiskiem uruchomieniowym .NET.
  \item Dostępna jest instancja serwera PostgreSQL; baza danych i jej schemat są tworzone automatycznie
        przy pierwszym uruchomieniu za pomocą migracji Entity Framework Core.
  \item Dane początkowe (przykładowy nauczyciel, uczniowie, grupy i testy) są wstawiane automatycznie,
        co umożliwia natychmiastowe przetestowanie aplikacji.
  \item Każdy nauczyciel pracuje na własnym zbiorze testów; grupy uczniów są zasobem współdzielonym.
\end{itemize}
